\par\noindent\textbf{(1)} 电报方程 \(\partial_zV=-L'\partial_tI\)、\(\partial_zI=-C'\partial_tV\) 导出波速 \(v=1/\sqrt{L'C'}\)，相位常数 \(\beta=\omega\sqrt{L'C'}\)。采用 \(e^{i\omega t}\) 约定，电压 \(V(z)=V_+e^{-i\beta z}+V_-e^{i\beta z}\)，电流 \(I(z)=Z_0^{-1}(V_+e^{-i\beta z}-V_-e^{i\beta z})\)。
\par\noindent\textbf{(2)} 令负载位于 \(z=0\)，边界 \(V(0)=Z_LI(0)\) 解得 \(\Gamma_L=V_-/V_+=(Z_L-Z_0)/(Z_L+Z_0)\)。对无损实负载，功率反射率 \(R=|\Gamma_L|^2\)，吸收份额 \(1-R\)，匹配 \(Z_L=Z_0\) 时无反射。
\par\noindent\textbf{(3)} 将行波解代入 \(Z_{\mathrm{in}}=V(-\ell)/I(-\ell)\)，利用正弦余弦恒等式化为 \(Z_{\mathrm{in}}=Z_0[Z_L+iZ_0\tan(\beta\ell)]/[Z_0+iZ_L\tan(\beta\ell)]\)。四分之一波长且负载为实正时 \(Z_{\mathrm{in}}=Z_0^2/Z_L\)，可实现阻抗变换。
\par\noindent\textbf{(4)} 对 \(|\Gamma_L|<1\)，驻波最大最小电压幅比分别为 \((1+|\Gamma_L|)/(1-|\Gamma_L|)\)。开路 \(Z_L\to\infty\) 时 \(\Gamma_L\to+1\)，短路 \(Z_L=0\) 时 \(\Gamma_L=-1\)，两者均全反射；若负载有增益或负阻，普通被动驻波比公式需重新界定。
